\chapter{Materiales}

En esta sección se definirán y precisarán todas las herramientas, métodos, algoritmos y equipos que fueron empleados y tomados para la realización del estudio.

\begin{table}[htbp]
\begin{center}
\begin{tabular}{|l||c|c|p{6cm}|} \hline
Nube & nPuntos & Dataset & Descripción \\ \hline
\textit{bildstein\_station1\_xyz\_intensity\_rgb} & 29.6 M & Semantic3D & Escaneo LiDAR terrestre y urbano desde una posición estática, muy denso. \\ \hline
\textit{sg27 station8 intensity rgb} & 73.9 M & Semantic3D & Escaneo LiDAR. \\ \hline
\textit{5145\_54340}                             & 12.1 M & DALES & LiDAR aéreo obtenido desde aeronaves. Está dividido en celdas disjuntas. \\ \hline
\textit{Lille\_0}                                & 10.0 M & Paris-Lille 3D & Escaneo . \\ \hline
\textit{Paris\_Luxembourg\_6}                    & 10.0 M & Paris-Lille 3D & Sensor LiDAR sobre un vehículo (MLS) que mapea dos calles en París. \\ \hline
\textit{PNOA\_2024\_PNR\_489-4672\_NPC01}        & 18.5 M & PNOA-II & Cobertura del territorio español mediante LiDAR aéreo. Tienen baja densidad. \\ \hline
\textit{Mar18_train} & 59.4 M & Hessigheim 3D & Dataset aéreo capturado mediante UAV (\textit{Unmanned Aerial Vehicle}) del pueblo alemán de Hessigheim\\ \hline
\textit{Mar18_test} & 59.4 M & Hessigheim 3D & Dataset aéreo capturado mediante UAV (\textit{Unmanned Aerial Vehicle}) del pueblo alemán de Hessigheim\\ \hline
\end{tabular}
\caption{Mapa de nubes utilizadas en las evaluaciones experimentales.}
\label{enlace2}
\end{center}
\end{table}